% Analytical core. Requires amsmath, amssymb and a natbib-compatible citation setup.
% Invoke \CoreDiscussion after the experimental/results sections.
\section{Introduction}
\label{sec:introduction}

Noise calibration affects both the strength of channel denoising and the
choice of a channel model in pilot-assisted orthogonal frequency-division
multiplexing (OFDM). When no independent noise measurement is available, a
receiver can estimate the noise variance from the portion of its pilot
observations left unexplained by an assumed channel subspace. This economical
approach couples calibration to the channel representation because unexplained
signal components enter the same residual energy as receiver noise. Increasing
the number of calibration blocks can therefore make a noise estimate more
repeatable without making it more accurate.

This distinction matters when a short integer-delay representation is used for
a channel with fractional path delays. The residual then contains channel
energy outside the selected subspace. Increasing the subspace dimension, however,
leaves fewer residual degrees of freedom for estimating noise. These effects
create a trade-off among channel approximation, calibration variance, and
observation cost. They also complicate the interpretation of downstream
performance. A change in channel-estimation error may reflect a
better representation, a different noise estimate, or a change in model order.

Established methods already address continuous channel parameters and unknown
statistics. Parametric OFDM receivers avoid an integer-delay restriction
\cite{hansen2018iterative}; variational line-spectrum estimation represents
continuous-parameter uncertainty and learns noise statistics
\cite{badiu2017valse}. Newtonized pursuit and its multiple-snapshot extension
provide further ways to refine continuous parameters
\cite{mamandipoor2016nomp,zhu2019mnomp}. These methods motivate a specific
diagnostic question: how much residual calibration error can additional
observations reduce, and how much requires a change in the signal model?

We quantify when residual calibration becomes dominated by channel leakage
and how temporal correlation changes the useful averaging window. The
projected channel covariance separates persistent bias, receiver-noise
fluctuations, and correlated leakage fluctuations. Their relative sizes give
a model-conditioned crossover between sampling-dominated and bias-dominated
calibration. Controlled OFDM comparisons then connect this distinction to
channel reconstruction: changing the channel representation and changing the
noise source address different parts of the error. A paired experiment
substitutes cyclic-prefix (CP) calibration within the same continuous-delay
receiver, using the existing waveform and preserving its noise dependence.
The resulting diagnostic identifies when additional averaging is useful and
when signal cancellation or a different channel representation is required.

\section{Related work}
\label{sec:related}

\subsection{Noise information and channel denoising}
Risk-based denoising offers an established alternative to assuming a known
channel covariance. Upadhya et al.\ use Stein's unbiased risk estimate to choose
linear and nonlinear OFDM channel estimators \cite{upadhya2016risk}. In their
accessible author preprint, the simulation setup estimates noise from 500 blank
OFDM carriers. This differs from estimating noise from a small residual
subspace of the pilots also used for channel estimation. Comparing these
settings requires distinguishing the denoising rule from the amount and source
of its noise information.

The CP provides noise information from samples already present in the waveform.
Cui and Tellambura estimate the power delay profile and noise variance from
the repeated signal structure; their Eq.~(15) gives a direct CP--tail
difference statistic \cite{cui2006pdpnoise}. Their simplified channel model
rounds delays to integers and assumes synchronization and blockwise channel
invariance. Here that established statistic serves as a calibration-source
control: a conservative delay bound selects a clean repetition interval, a
finite-burst waveform check verifies fractional-delay cancellation, and paired
receiver comparisons retain the dependence between CP and FFT noise. This
separates the source of noise information from the denoising rule that uses it.

\subsection{Continuous parameters and multiple observations}
Continuous-delay modeling has a direct counterpart in line-spectrum
estimation. Newtonized orthogonal matching pursuit (NOMP) alternates detection
with refinement of sinusoidal parameters and uses a noise-dependent stopping
criterion \cite{mamandipoor2016nomp}. Its multiple-snapshot extension, MNOMP,
uses observations with shared frequencies and snapshot-dependent amplitudes
\cite{zhu2019mnomp}. For OFDM, Gao et al.\ exploit common delay structure and
spatial and temporal correlations \cite{gao2014superresolution}. These
approaches use shared channel structure across observations, whereas residual
energy pooling shares a scalar noise statistic. The distinction matters when
comparing observation budgets and attributing gains to calibration.

Unknown noise and continuous parameters can also be inferred jointly. VALSE
maintains approximate frequency distributions, selects a model order, and
updates the noise and gain statistics \cite{badiu2017valse}. Hansen et al.\
incorporate parametric channel estimation into an iterative OFDM receiver that
exchanges soft information with a decoder and represents dense channels through
a virtual channel impulse response \cite{hansen2018iterative}. These methods
establish joint continuous-parameter and noise inference; the calibration
question here concerns the consequences of retaining a fixed residual subspace.

Subsequent variational formulations extend uncertainty-aware inference to
multidimensional spectra \cite{zhang2020mdvalse} and noise variances that vary
across sensors or snapshots \cite{zhang2021heteroscedastic}. For underwater
acoustic OFDM, Wan et al.\ jointly estimate frequency offset, continuous-delay
channels, and noise statistics while exchanging information with a decoder;
their initialization also uses pilot-only VALSE \cite{wan2022jccdvalse}.
These developments reinforce the distinction between learning nuisance
parameters within a channel model and examining the consequences of fixing an
imperfect calibration subspace.

Table~\ref{tab:priorwork} organizes these precedents by the information they
use. Our contribution is the quantitative connection between a fixed
subspace's leakage, the averaging window, and the reconstruction gain obtained
by changing the calibration source. Standard Gaussian quadratic-form identities
provide the analytical basis. A specified orthogonal matching pursuit
(OMP)/variable-projection receiver provides the controlled intervention, and
a separate paper-equation VALSE implementation examines joint inference in
the pilot domain.

\begin{table}[t]
\centering
\caption{Closest prior approaches and the calibration question examined here.
The rows distinguish information use, rather than numerical performance.}
\label{tab:priorwork}
\small
\renewcommand{\arraystretch}{1.12}
\begin{tabular}{@{}>{\raggedright\arraybackslash}p{0.23\linewidth}>{\raggedright\arraybackslash}p{0.33\linewidth}>{\raggedright\arraybackslash}p{0.38\linewidth}@{}}
\toprule
Approach & Established mechanism & Relation to this study \\
\midrule
Risk-based denoising \cite{upadhya2016risk} & SURE selects channel denoisers;
the author preprint uses 500 blank carriers for noise estimation. & Separates
the denoiser from the quantity and source of noise observations. \\
CP calibration \cite{cui2006pdpnoise} & Direct CP--tail differences cancel
repeated signal samples. & Supplies the established statistic for a paired
calibration-source substitution with fractional delays. \\
NOMP/MNOMP and common-delay OFDM \cite{mamandipoor2016nomp,zhu2019mnomp,gao2014superresolution}
& Continuous parameters and shared geometry improve channel fitting. &
Distinguishes sharing channel structure from pooling residual energies. \\
VALSE and iterative OFDM receivers \cite{badiu2017valse,hansen2018iterative,wan2022jccdvalse}
& Joint inference learns continuous parameters, model order, and noise;
iterative receivers also use decoder information. & Motivates the pilot-domain
joint-inference check and fixes the information boundary of the receiver controls. \\
\bottomrule
\end{tabular}
\end{table}

\section{System model}
\label{sec:system}

\subsection{OFDM observations and delay convention}
Consider a single-antenna OFDM link with an even number $N$ of subcarriers. Let
$q_k=k$ for $0\leq k<N/2$ and $q_k=k-N$ otherwise, and set $f_k=q_k/N$ in cycles
per base-rate sample. The frequency response in block $t$ is
\begin{equation}
 H_{t,k}=\sum_{j=1}^{J}g_{t,j}\exp(-\mathrm{i}2\pi f_k\tau_j),
 \qquad k=0,\ldots,N-1,
 \label{eq:channel}
\end{equation}
where $\tau_j$ is measured in sample periods and delays remain fixed over a
calibration window. The path gains are zero-mean proper complex Gaussian with
covariance $\mathbf P=\operatorname{diag}(p_1,\ldots,p_J)$ and
$\sum_jp_j=1$. This normalization applies to the ensemble; individual channel
realizations are not rescaled.

There are $M$ equally spaced pilots of value one at indices $k_m=mD$, where
$D=N/M$. After division by the known pilots, the observations satisfy
\begin{equation}
 \mathbf y_t=\mathbf B\mathbf g_t+\mathbf n_t,
 \quad B_{mj}=\exp(-\mathrm{i}2\pi f_{k_m}\tau_j),
 \quad \mathbf n_t\sim\mathcal{CN}(\mathbf0,v\mathbf I_M).
 \label{eq:pilotmodel}
\end{equation}
Noise is independent across blocks and of all gains. The scalar
$v$ is the variance per complex pilot observation and is constant within the
window. Under the normalization above, the nominal subcarrier SNR is $1/v$,
excluding CP and pilot overhead. True delays, powers, and $v$ are used for data
generation and explicitly identified oracle references; practical estimators
do not receive them.

The numerical setting uses $N=64$, $M=16$, and a nominal integer-delay window
of length $L=12$. For fractional delays, replacing a negative frequency by its
unsigned FFT index changes the continuous-delay phase, so the signed-frequency
convention matters. Pilot delay atoms have a period of $M$ samples in this
setting, and the predefined search interval $[0,12]$ lies within one such
period. All generators and estimators use the same convention.

\subsection{Waveform assumptions and calibration windows}
For a useful symbol and its CP, the ideal transmitted waveform is
\begin{equation}
 s_t(u)=\frac{1}{\sqrt N}\sum_{k=0}^{N-1}X_{t,k}
             \exp(\mathrm{i}2\pi f_ku),
 \qquad -N_{\rm CP}\leq u<N.
 \label{eq:waveform}
\end{equation}
Adjacent CP-extended symbols occupy separate rectangular intervals. With ideal
delayed paths, exact synchronization, constant gains during the received
symbol, and $N_{\rm CP}\geq\max_j\tau_j$, CP removal and a unitary FFT give
the multiplicative response in \eqref{eq:channel}. To check finite-CP
consistency for ideal paths, the waveform audit applies linear delays to a
finite oversampled burst rather than a circular shift to each symbol
separately. The rectangular symbol boundaries do not form a globally bandlimited
analog waveform. Transmit/receive filtering, frequency offsets, Doppler, and
radio-frequency impairments are outside this model.

The main calibration experiment holds path geometry and noise variance fixed
while drawing independent gains for each block; a separate analysis allows
stationary temporal correlation. Independent fading blocks are not taken to be
consecutive OFDM symbols, because the observation duration needed to obtain
them depends on channel coherence. Only the first block is the target of
channel estimation and data detection. The other
blocks supply scalar calibration statistics unless a method is explicitly
defined to fit shared channel parameters.

\section{Analytical characterization of residual calibration}
\label{sec:analysis}

\subsection{Fixed subspaces and conditional moments}
Let $\mathbf A\in\mathbb C^{M\times L}$ contain the pilot responses of integer
delays $0,\ldots,L-1$, with $L<M$. For the comb pilots,
$\mathbf A^{\mathrm H}\mathbf A=M\mathbf I_L$. Define
\begin{equation}
 \mathbf Q=\mathbf I_M-\frac{\mathbf A\mathbf A^{\mathrm H}}{M},
 \quad d=M-L,\quad
 \mathbf z_t=\frac{\mathbf A^{\mathrm H}\mathbf y_t}{M},
 \quad R_t=\mathbf y_t^{\mathrm H}\mathbf Q\mathbf y_t.
 \label{eq:projection}
\end{equation}
$\mathbf Q$ is a fixed Hermitian projector with rank $d$. For a general
full-rank pilot basis, the projector is instead
$\mathbf I_M-\mathbf A(\mathbf A^{\mathrm H}\mathbf A)^{-1}\mathbf A^{\mathrm H}$;
the residual identities below remain valid. Pooling $T$ blocks gives
\begin{equation}
 \widehat v_T=\frac{1}{Td}\sum_{t=1}^{T}R_t.
 \label{eq:vhat}
\end{equation}

First condition on the noiseless observations
$\boldsymbol\mu_t=\mathbf B\mathbf g_t$. For proper complex Gaussian noise,
expanding $R_t$ into its deterministic, linear-noise, and quadratic-noise
terms yields
\begin{align}
 \mathbb E[R_t\mid\boldsymbol\mu_t]
   &=vd+\|\mathbf Q\boldsymbol\mu_t\|^2,\label{eq:conditionalmean}\\
 \operatorname{Var}(R_t\mid\boldsymbol\mu_t)
   &=v^2d+2v\|\mathbf Q\boldsymbol\mu_t\|^2.
 \label{eq:conditionalvariance}
\end{align}
The linear term has variance $2v\|\mathbf Q\boldsymbol\mu_t\|^2$;
the quadratic term has variance $v^2\operatorname{tr}(\mathbf Q^2)=v^2d$.
Their covariance is zero because centered Gaussian third moments vanish.
Since receiver noise is independent across blocks,
\begin{align}
 \mathbb E[\widehat v_T\mid\{\boldsymbol\mu_t\}]
  &=v+\frac{\sum_t\|\mathbf Q\boldsymbol\mu_t\|^2}{Td},
  \label{eq:poolconditionalmean}\\
 \operatorname{Var}(\widehat v_T\mid\{\boldsymbol\mu_t\})
  &=\frac{Tv^2d+2v\sum_t\|\mathbf Q\boldsymbol\mu_t\|^2}{T^2d^2}.
  \label{eq:poolconditionalvariance}
\end{align}
If the signal belongs to the assumed subspace, all projected means vanish.
Then $\sum_tR_t/v\sim\operatorname{Gamma}(Td,1)$, so
$\mathbb E[\widehat v_T]=v$ and
$\operatorname{Var}(\widehat v_T)=v^2/(Td)$. These statements remain true
for correlated gains when the subspace is correct, as the residual contains
receiver noise alone. Increasing $L$ reduces $d$, so a wider delay window
increases calibration variance even without mismatch.

\subsection{Ensemble leakage bias and Gaussian variance}
To average over channel realizations, define
\begin{equation}
 \mathbf C_s=\mathbf B\mathbf P\mathbf B^{\mathrm H},\qquad
 \mathbf S=\mathbf Q\mathbf C_s\mathbf Q,\qquad
 \mathbf C=\mathbf S+v\mathbf Q,\qquad
 b=\frac{\operatorname{tr}(\mathbf S)}{d}.
 \label{eq:leakagecovariance}
\end{equation}
Here $\mathbf S$ is the covariance of the signal remaining after projection;
$b\geq0$ is its residual energy per retained complex dimension. Since
$\mathbf Q^2=\mathbf Q$,
$\operatorname{tr}(\mathbf S)=\operatorname{tr}(\mathbf Q\mathbf C_s)$.
The projected observation is proper Gaussian with covariance $\mathbf C$,
possibly singular in the ambient pilot space. Diagonalizing this covariance
expresses its squared norm as a weighted sum of independent unit-mean
exponential variables, giving
\begin{equation}
 \mathbb E[R_t]=\operatorname{tr}(\mathbf C),\qquad
 \operatorname{Var}(R_t)=\operatorname{tr}(\mathbf C^2).
 \label{eq:singlemoments}
\end{equation}
The mean identity requires the stated second moments, whereas the variance
identity also uses proper Gaussian gains. For independent blocks,
\begin{align}
 \mathbb E[\widehat v_T]&=v+b,\label{eq:bias}\\
 \operatorname{Var}(\widehat v_T)
 &=\frac{\operatorname{tr}(\mathbf S^2)
          +2v\operatorname{tr}(\mathbf S)+v^2d}{Td^2}.
 \label{eq:independentvariance}
\end{align}
Pooling reduces variance around $v+b$ but leaves $b$ unchanged. For fixed
channel power and geometry, the relative bias $b/v$ increases with SNR.
These quantities separate thermal-noise estimation error from residual
signal-model error and can be computed before Monte Carlo simulation.

\subsection{Correlated calibration blocks}
Consider stationary Gaussian gains with a real autoregressive coefficient
$0\leq\rho\leq1$:
\begin{equation}
 \mathbf g_{t+1}=\rho\mathbf g_t+\sqrt{1-\rho^2}\,\mathbf w_t,
 \quad \mathbf g_1,\mathbf w_t\sim\mathcal{CN}(\mathbf0,\mathbf P),
 \label{eq:ar1}
\end{equation}
where the innovations are mutually independent and independent of receiver
noise. Thus
$\mathbb E[\mathbf g_t\mathbf g_u^{\mathrm H}]=\rho^{|t-u|}\mathbf P$.
For $t\ne u$, the residual cross-covariance is
$\rho^{|t-u|}\mathbf S$. The proper Gaussian fourth-moment identity gives
\begin{equation}
 \operatorname{Cov}(R_t,R_u)
   =\rho^{2|t-u|}\operatorname{tr}(\mathbf S^2).
 \label{eq:crosscovariance}
\end{equation}
Pairing the two cross-block factors in the expansion of
$\mathbb E[R_tR_u]$ produces the trace of the cross-covariance times its
Hermitian transpose; the within-block pairing cancels the product of means.
Summing all diagonal and off-diagonal terms yields
\begin{equation}
 \operatorname{Var}(\widehat v_T)=
 \frac{T\operatorname{tr}(\mathbf C^2)
 +2\operatorname{tr}(\mathbf S^2)
   \sum_{h=1}^{T-1}(T-h)\rho^{2h}}{T^2d^2}.
 \label{eq:correlatedvariance}
\end{equation}
The expectation remains \eqref{eq:bias}. To show which part of the variance
is affected by correlation, set
$a_T=T+2\sum_{h=1}^{T-1}(T-h)\rho^{2h}$. Then
\begin{equation}
 \operatorname{Var}(\widehat v_T)=
 \frac{v^2d+2v\operatorname{tr}(\mathbf S)}{Td^2}
 +\frac{a_T\operatorname{tr}(\mathbf S^2)}{T^2d^2}.
 \label{eq:varianceparts}
\end{equation}
The effective sample size $T^2/a_T$ describes the channel-leakage term,
not the entire noise estimate. For fixed $\rho<1$,
$a_T/T\to(1+\rho^2)/(1-\rho^2)$, and the variance tends to zero while the
estimate concentrates around $v+b$.

When $\rho=1$, all blocks share one gain realization, $a_T=T^2$, and the
ensemble variance instead approaches
$\operatorname{tr}(\mathbf S^2)/d^2$. Conditional on that realization,
averaging still removes receiver-noise fluctuations, and the estimate converges
to $v+\|\mathbf Q\mathbf B\mathbf g_1\|^2/d$. An ensemble variance floor
therefore does not mean that receiver-noise averaging has stopped working.
In all cases, the mean-square calibration error is
\begin{equation}
 \mathbb E[(\widehat v_T-v)^2]
   =b^2+\operatorname{Var}(\widehat v_T),
 \label{eq:calibrationmse}
\end{equation}
showing the persistent contribution of bias.

\subsection{A crossover for the value of additional averaging}
\label{sec:crossover}
For independent blocks and $b>0$, write
\begin{equation}
 A=\frac{\operatorname{tr}(\mathbf S^2)+2v\operatorname{tr}(\mathbf S)+v^2d}{d^2},
 \qquad T_\star=\frac{A}{b^2}.
 \label{eq:crossover}
\end{equation}
Equation~\eqref{eq:calibrationmse} then becomes $b^2+A/T$.
The fraction attributable to squared bias is
$\phi_T=b^2/(b^2+A/T)=T/(T+T_\star)$.
For $T\geq T_\star$, leakage bias accounts for at least half the mean-square
calibration error; further pooling reduces its variance contribution. When $b=0$,
the error follows $v^2/(Td)$ throughout. Thus $T_\star$ locates a change in
the dominant error source, conditional on the channel covariance and noise
level. It is a design diagnostic rather than an estimated stopping rule.

For correlated gains, the same comparison uses
\eqref{eq:varianceparts} at the chosen $(T,\rho)$. Correlation changes the
leakage-fluctuation term while the receiver-noise terms retain their $1/T$
scaling. At $\rho=1$, both the squared bias and the random channel-specific
offset persist. These expressions distinguish three responses to a
calibration gap: average while sampling variance dominates, change the signal
subspace to alter leakage, or use signal-cancelling observations to estimate
the thermal variance directly.

\subsection{Dependence and fitted-model boundaries}
In the correctly specified model, the target-block channel statistic and its
residual noise are orthogonal. Under mismatch,
\begin{equation}
 \mathbb E[\mathbf z_t(\mathbf Q\mathbf y_t)^{\mathrm H}]
       =\frac{\mathbf A^{\mathrm H}\mathbf C_s\mathbf Q}{M}
 \label{eq:statisticdependence}
\end{equation}
need not vanish. Although their receiver-noise components are independent,
their signal components can couple channel estimation to calibration. Moving
from $T=1$ to $T>1$ changes both the dispersion of the noise estimate and its
dependence on the target block, so a downstream error difference cannot be
attributed entirely to variance reduction.

The moment identities describe a projector fixed independently of the observed
data. For fitted delays or support, the projector is data-dependent and its
residual distribution includes the fitting and selection process. Counting
each fitted real delay as an additional real parameter gives a local
degrees-of-freedom approximation whose accuracy depends on order selection
and path resolution. Jointly fitted models therefore need their own
uncertainty treatment, as in VALSE \cite{badiu2017valse}.

\subsection{A cyclic-prefix calibration control}
\label{sec:cpcontrol}
The CP--tail difference estimator of Cui and Tellambura
\cite[Eq.~(15)]{cui2006pdpnoise} cancels the repeated signal before measuring
noise energy. We use this established construction with a conservative memory
bound to isolate calibration-source effects. With $N_{\rm CP}=16$ and the
known delay bound $\max_j\tau_j\leq12$, the last $K_{\rm CP}=4$ base-rate CP
samples and their counterparts at the end of the useful symbol have equal
noiseless received values. For received time samples $r_t[u]$, define
\begin{equation}
 \widehat v_{{\rm CP},T}
 =\frac{1}{2TK_{\rm CP}}
   \sum_{t=1}^{T}\sum_{a=1}^{K_{\rm CP}}
       |r_t[-a]-r_t[N-a]|^2.
 \label{eq:cpnoise}
\end{equation}
Independent base-rate receiver noise makes each difference
$\mathcal{CN}(0,2v)$, giving exact moments
\begin{equation}
 \mathbb E[\widehat v_{{\rm CP},T}]=v,\qquad
 \operatorname{Var}(\widehat v_{{\rm CP},T})
       =\frac{v^2}{TK_{\rm CP}}.
 \label{eq:cpmoments}
\end{equation}
Using existing CP samples and the declared delay bound, this control adds no
transmitted pilots. It requires synchronization, an adequate channel memory
margin, and unchanged gains over each paired interval.

The paired noise differences must remain coupled to the pilot
observations. If $\mathbf F$ is the unitary FFT matrix, $\mathcal I$ indexes the
paired useful-symbol samples, and
$\mathbf e=(\mathbf n_{\rm CP}-\mathbf n_{\mathcal I})/\sqrt2$, then
\begin{equation}
 \mathbb E[\mathbf n_{\mathcal P}\mathbf e^{\mathrm H}]
       =-\frac{v}{\sqrt2}\mathbf F_{\mathcal P,\mathcal I},
 \label{eq:cpdependence}
\end{equation}
where $\mathbf n_{\mathcal P}$ denotes the FFT noise on the pilots.
A separate independent calibration sample with the same marginal distribution
would discard this dependence. We therefore form CP differences and FFT
observations from the same time-domain noise in the receiver comparison.

\newcommand{\CoreDiscussion}{%
\section{Discussion}
\label{sec:discussion}

Calibration-source substitution separates a recoverable noise-calibration
gap from the error imposed by the channel representation. In the specified
continuous-delay receiver, CP calibration recovers most of the performance
available with known noise. Applying the same statistic to integer EB yields
a smaller gain because its representation still excludes fractional-delay
energy. This paired contrast links the projected leakage identified by the
analysis to a consequential receiver decision: an inflated noise estimate
changes the penalty for retaining weak components.

The bias--variance crossover makes the averaging trade-off explicit. Once
$b^2$ exceeds the variance at the selected window, additional residual samples
mainly sharpen an estimate centered above the thermal variance. Temporal
correlation extends the window needed to suppress leakage fluctuations; with
constant gains, the estimate stabilizes around a channel-specific offset.
Changing the integer window alters both leaked energy and residual dimension,
so its calibration effect is determined by their ratio. These quantities
provide a common basis for comparing more observations, a different subspace,
and a signal-cancelling noise source.

Information budgets determine how these alternatives can be used. The CP
control uses additional received samples from the existing waveform and a
known memory bound. Independent probes supply separate signal-free
observations. Shared-delay methods additionally exploit channel structure
across blocks \cite{zhu2019mnomp,gao2014superresolution}. Accordingly, an
observation window is characterized by its pilots, waiting duration,
stationarity assumptions, and shared parameters. Within the present paired
CP comparison, the channel data and receiver settings are held fixed while
the calibration source changes; the resulting effect includes model-order
selection through Eq.~\eqref{eq:bic}. The grid and refined controls also
change parameter penalties and serve as model comparisons.

The scope is fixed-subspace calibration under the stated synchronized OFDM
model and its explicitly specified receiver controls. The finite-burst check
validates fractional delays for ideal paths; filtering requires the
corresponding effective channel and noise covariance. CP cancellation requires
a clean repeated interval, sufficient effective memory margin, and unchanged
gains across each pair. The dense-channel comparison concerns the chosen
six-path control and its selection and optimization settings; parametric
receivers with richer representations
can also model dense channels \cite{hansen2018iterative}. The separate VALSE
experiment examines pilot-domain joint inference, with its own support and
observation budget. These boundaries connect each numerical result to the
calibration or representation choice that generated it.
}
